\documentclass[11pt,a4paper]{article}

% --- Packages ---
\usepackage[margin=.75in]{geometry}
\usepackage[utf8]{inputenc}
\usepackage[T1]{fontenc}
\usepackage{titlesec}
\usepackage{enumitem}
\usepackage{hyperref}
\usepackage{cite}

% Section formatting - clean and academic
\titleformat{\section}{\large\bfseries}{\thesection}{1em}{}[\titlerule]
\titlespacing*{\section}{0pt}{3.5ex plus 1ex minus .2ex}{2.3ex plus .2ex}

% --- Metadata ---
\title{
    \vspace{-1in}
    \large \textbf{Research Statement}\\
    \huge \textbf{Towards hardware-accelerated data-planes inference}\\
    \large \textit{}
}
\author{\textbf{Giuseppe Capasso} \\ \href{mailto:capassog97@gmail.com}{capassog97@gmail.com} | \href{https://giuseppe.capass.org}{giuseppe.capass.org} | \href{https://github.com/alarmfox}{github.com/alarmfox}}
\date{}

\begin{document}

\maketitle
I have a hybrid background in System Architectures, Computer Networks and Cybersecurity. During my research internship, I focused on VM-based Trusted Execution Environments (TEEs) like Intel TDX and AMD SEV, designing a trusted firmware and hypervisor for RISC-V architectures to support isolated, bare-metal, virtualized workloads adhering to the CoVE specification. I also contributed to Simply-V, a RISC-V soft SoC for Xilinx FPGAs, building its software stack from scratch and porting FreeRTOS and OpenSBI. My specialization in TEEs began during my Master's, where I developed a Rust CLI tool leveraging eBPF, perf, and Intel RAPL to benchmark Intel SGX enclaves for performance and power.

In parallel, I focused on high-performance networking architectures implementing low-latency streaming applications using the QUIC protocol, alongside exploring cloud network infrastructure. Finally, for the last five years, I have been an instructor and mentor at Cisco DTLab in Naples holding courses about CCNA, Cybersecurity and Linux, leading workshops on containers and infrastructure automation.

Driven by these combined interests, my research vision focuses on the intersection of hardware-accelerated network security and modern cloud architectures. Specifically, I am deeply interested in advancing Software-Defined Networking (SDN) and Network Functions Virtualization (NFV) paradigms, including highly optimized load balancers and hardware-accelerated Network Intrusion Detection Systems (NIDS). A critical frontier in this domain is high-performance traffic characterization, which is essential for dynamic routing optimization and managing the latency-sensitive workloads generated by emerging technologies like Large Language Models (LLMs).

However, performing line-rate inference and feature extraction at Layer 3 and Layer 7 may introduce privacy risks. Analyzing this traffic exposes highly sensitive user data (such as DNS queries and proprietary payloads) to the underlying infrastructure. To overcome these privacy barriers, my research aims to integrate Confidential Computing directly into the network data plane. By leveraging hardware-assisted TEEs, we can perform high-performance, on-chip threat inference while cryptographically isolating user data. Ultimately, this challenge perfectly aligns with my passion for low-level system architectures and network security, offering a clear path to solving critical, real-world vulnerabilities in next-generation cloud environments.

\end{document}
